\documentclass[11pt,letterpaper]{article}
\usepackage[T1]{fontenc}
\usepackage[utf8]{inputenc}
\usepackage{lmodern}
\usepackage[margin=1in]{geometry}
\usepackage{amsmath}
\usepackage{booktabs}
\usepackage{longtable}
\usepackage{enumitem}
\usepackage{listings}
\usepackage{xcolor}
\usepackage[hidelinks]{hyperref}

\definecolor{backcolour}{rgb}{0.96,0.96,0.94}
\definecolor{cautionred}{rgb}{0.55,0,0}

% Listings are configured so that anything inside one can be selected from the
% PDF and pasted straight into a terminal:
%   numbers=none        no line numbers to catch in the selection
%   columns=fullflexible no padding spaces inserted between characters
%   upquote=true        straight quotes, not typographic ones
%   breaklines=false    a wrapped line would paste with a spurious newline,
%                       so every line below is kept short enough to fit
\lstset{
    basicstyle=\ttfamily\small,
    backgroundcolor=\color{backcolour},
    numbers=none,
    frame=single,
    framesep=5pt,
    columns=fullflexible,
    keepspaces=true,
    upquote=true,
    showstringspaces=false,
    breaklines=false,
    xleftmargin=0pt,
    aboveskip=8pt,
    belowskip=8pt
}

\newenvironment{caution}[1]{%
    \par\addvspace{8pt}%
    \noindent{\color{cautionred}\rule{\linewidth}{0.8pt}}\par\vspace{2pt}
    \noindent\textbf{\color{cautionred}Caution: #1}\par\vspace{2pt}
}{%
    \par\vspace{2pt}\noindent{\color{cautionred}\rule{\linewidth}{0.4pt}}%
    \par\addvspace{8pt}%
}

\title{Arctos CAN Interface \\ \large How the arm's motors and CAN bus actually work}
\author{Consolidated from project debugging logs, source inspection and bench testing}
\date{September 2026}

\begin{document}
\maketitle
\tableofcontents
\newpage

\section{About these documents}

This is the reference document: how the bus works, how the protocol is framed, how
to turn encoder counts into joint angles, and what has actually been verified
against hardware rather than assumed. It is deliberately not a step-by-step
procedure.

There are four documents in this set and they do not overlap:

\begin{itemize}[leftmargin=*]
    \item \textbf{Arctos\_CAN\_Interface} --- this one. Background and reference.
    \item \textbf{Nema\_17\_MKS42D\_SOP} --- bringing up a NEMA 17 joint on an MKS
          SERVO42D. Joints A, B and C.
    \item \textbf{Nema\_23\_MKS57D\_SOP} --- bringing up a NEMA 23 joint on an MKS
          SERVO57D. Joint X. Largely unverified; see that document's opening.
    \item \textbf{BASIC\_CAN\_COMMANDS} --- terminal commands only. Setting the
          interface up, talking to a motor, shutting down to unplug.
\end{itemize}

This set replaces \texttt{arctos\_motor\_guide}, \texttt{sop\_nema17},
\texttt{sop\_nema23}, \texttt{arctos\_CAN\_ROS2\_SOP\_updated} and
\texttt{arctos\_motor\_can\_cheat\_sheet\_updated}. Section~\ref{sec:corrections}
lists the things that were \emph{wrong} in those files, so nothing is silently
re-derived from a stale copy.

\section{Core concepts}

\subsection{CAN bus, in one paragraph}
CAN is a differential two-wire multi-drop bus. Every device shares the same pair,
CAN\_H and CAN\_L, plus a common ground, and each device is addressed by an ID.
Every node must agree on the bitrate; this arm runs at \textbf{500 kbit/s}. Because
all nodes share one pair, two devices configured with the same CAN ID will collide
and neither will behave correctly.

\subsection{Gearboxes and closed-loop steppers, in one paragraph}
Each joint is a stepper motor driving a reduction gearbox. The reduction multiplies
torque and divides both speed and angular error: through a 67.82:1 gearbox, one
motor revolution is only about 5.31\textdegree{} at the joint. The MKS controllers
are closed-loop --- a magnetic encoder on the motor shaft reports true position, so
the controller can tell that it has fallen behind. Two consequences matter
constantly in practice. The encoder measures the \textbf{motor shaft, not the joint
output}. And a stalled stepper converts its full phase current into heat, with no
motion to carry it away.

\section{Hardware}

\subsection{Joint, motor and CAN ID map}

\begin{center}
\begin{tabular}{lllll}
\toprule
\textbf{Joint} & \textbf{Role} & \textbf{CAN ID} & \textbf{Motor / driver} & \textbf{Gear ratio} \\ \midrule
X & Base and waist rotation & \texttt{0x01} & NEMA 23 / SERVO57D & 13.5:1  \\
Y & not confirmed           & \texttt{0x02} & ---                & 150.0:1 \\
Z & not confirmed           & \texttt{0x03} & ---                & 150.0:1 \\
A & Elbow                   & \texttt{0x04} & NEMA 17 / SERVO42D & 48.0:1  \\
B & Wrist rotation          & \texttt{0x05} & NEMA 17 / SERVO42D & 67.82:1 \\
C & Wrist joint             & \texttt{0x06} & NEMA 17 / SERVO42D & 67.82:1 \\
\bottomrule
\end{tabular}
\end{center}

\begin{caution}{A driver can be found on the wrong CAN ID --- fix the panel, not the map}
The table above is the \emph{intended} assignment and matches
\texttt{arctos\_controller.yaml}. A controller's panel can hold a different value.
On 2026-09-02 a scan with only Joint B connected returned a single responder at
\texttt{0x06}, which is Joint C's ID: that driver's panel had simply been saved
wrong. The work done that day was on Joint B regardless of the ID it answered on.

When a scan disagrees with this table, that is evidence the \textbf{panel needs
reprogramming}, not evidence the map is wrong. Reprogram the driver rather than
editing the map, and never drive a joint until the scan and the map agree. Joint
A's \texttt{0x04} and its 48:1 ratio are empirically verified; the rest are config
values.
\end{caution}

\begin{caution}{Always confirm a CAN ID against real hardware}
Run the read-only scan before trusting any table, including the one above. It sends
only the \texttt{0x31} query and never enables coils:
\begin{lstlisting}
cd ~/arctos_ws/src/arctos_can_control/arctos_can_control
python3 can_id_scan.py --channel can0 --ids 1-16
\end{lstlisting}
\end{caution}

\subsection{Motor specifications}

\begin{center}
\begin{tabular}{lll}
\toprule
\textbf{Parameter} & \textbf{NEMA 17 / SERVO42D} & \textbf{NEMA 23 / SERVO57D} \\ \midrule
Holding torque & 24~N\textperiodcentered cm & 1.8~N\textperiodcentered m \\
Rated current  & 1.2~A & 2.8~A \\
Shaft diameter & 5~mm  & 6.35~mm \\
Step angle     & 1.8\textdegree{}, 200 full steps per rev & 1.8\textdegree{}, 200 full steps per rev \\
Microstepping  & 16 times, 3200 pulses per rev & 16 times, 3200 pulses per rev \\
Encoder        & 16384 counts per rev, 14-bit & 16384 counts per rev, 14-bit \\
Joints         & A, B, C & X \\
\bottomrule
\end{tabular}
\end{center}

\begin{caution}{Do not confuse motor steps, microstep pulses and encoder counts}
Three different resolutions are in play, and mixing them is the single most common
arithmetic error in this project.
\begin{itemize}[leftmargin=*]
    \item \textbf{200} full steps per motor revolution, the physical stepping
          resolution.
    \item \textbf{3200} microstep pulses per motor revolution, the scale that
          \texttt{0xFD} position commands are expressed in.
    \item \textbf{16384} encoder counts per motor revolution, the scale that
          \texttt{0x31} reports in.
\end{itemize}
One microstep pulse is $16384/3200 = 5.12$ encoder counts. Commanded motion and
measured motion are therefore never in the same units.
\end{caution}

\subsection{CANable 2 adapter and wiring}

\begin{itemize}[leftmargin=*]
    \item USB VID:PID \texttt{16d0:117e}.
    \item It contains its own CAN controller. Do \textbf{not} add an external MCP2515
          module in series with it.
    \item Only three wires are needed between adapter and controller: CAN\_H, CAN\_L
          and GND. Do \emph{not} connect the CANable's power lines to the motor
          controller.
\end{itemize}

\begin{lstlisting}
Ubuntu host -> USB -> CANable 2 -> CAN_H / CAN_L / GND
            -> MKS controller (SERVO42D or SERVO57D) -> joint motor
\end{lstlisting}

\subsection{Bus termination check, power off}

With everything powered off, measure resistance between CAN\_H and CAN\_L.

\begin{center}
\begin{tabular}{ll}
\toprule
\textbf{Reading} & \textbf{Meaning} \\ \midrule
About 60~$\Omega$  & Correct. Two 120~$\Omega$ terminators in parallel \\
About 120~$\Omega$ & Only one terminator present \\
Open circuit       & Wiring is disconnected \\
Very low           & Possible short \\
\bottomrule
\end{tabular}
\end{center}

\section{The CAN interface on Ubuntu}
\label{sec:can-interface}

The CANable presents as a serial device, and \texttt{slcand} bridges it to a Linux
network interface. This has to be rebuilt every time the adapter is plugged in or
re-enumerates. The commands themselves are in \texttt{BASIC\_CAN\_COMMANDS}; what
follows is why they are what they are.

\begin{caution}{The \texttt{slcand} bitrate flag is a coded index, not a number}
\texttt{-s0} is 10k, \texttt{-s1} 20k, \texttt{-s2} 50k, \texttt{-s3} 100k,
\texttt{-s4} 125k, \texttt{-s5} 250k, \textbf{\texttt{-s6} 500k}, \texttt{-s7} 800k,
\texttt{-s8} 1M. This arm is 500 kbit/s, so \texttt{-s6} is correct. Earlier scripts
in this project used \texttt{-s3} and silently attached at 100 kbit/s, which
presents as a completely dead bus.
\end{caution}

\begin{caution}{On slcan the bitrate is set at attach time and \texttt{ip} cannot show it}
Do not use \texttt{ip link set can0 up type can bitrate 500000}. That netlink
attribute is for native CAN controller drivers and is not supported on slcan. As a
result \texttt{ip -details link show can0} reports \texttt{bitrate 0} and
\texttt{clock 0} on a perfectly healthy CANable. This is normal and is not a fault
to chase. An older version of this documentation told the reader to verify by
looking for \texttt{bitrate 500000}; that check can never pass on this hardware.
\end{caution}

\begin{caution}{Kernel CAN error counters are meaningless on slcan}
The re-started, bus-errors, arbit-lost, error-warn, error-pass and bus-off counters
in \texttt{ip -details -statistics link show can0} are not populated by the slcan
driver. They read zero regardless of what the bus is doing. Do not use them as
evidence that a bus is healthy, and note that the ``total silence with zero bus
errors implies no CAN transceiver'' heuristic in Section~\ref{sec:rs485} needs a
native CAN controller to mean anything.
\end{caution}

\section{Controller panel configuration}
\label{sec:panel}

Confirm these on the controller's own display before attempting CAN communication.

\begin{center}
\begin{tabular}{ll}
\toprule
\textbf{Setting} & \textbf{Value} \\ \midrule
CAN ID   & the joint's ID from the map above \\
CAN Rate & \texttt{500} \\
CAN Response, shown as CANRSP & \texttt{ENABLE} \\
CAN CRC  & \texttt{SUM-8} \\
Working current & at or below the motor's rated current \\
Working mode & see the caution below \\
\bottomrule
\end{tabular}
\end{center}

Save and power-cycle the controller after any change. The CAN ID in particular must
be committed, not merely displayed.

\begin{caution}{Set the working current to the motor's rating, not the panel default}
A controller in this project was found storing a 3000~mA default against a motor
rated 1.2~A, roughly two and a half times over. Check this on every new controller.
\end{caution}

\begin{caution}{A working-mode mismatch can imitate a communication fault}
The ROS 2 driver defines six modes and defaults to \texttt{CR\_vFOC}, which is 2:
\begin{lstlisting}
CR_OPEN = 0, CR_CLOSE = 1, CR_vFOC = 2,
SR_OPEN = 3, SR_CLOSE = 4, SR_vFOC = 5
\end{lstlisting}
If the panel's mode disagrees with what a motion command assumes, the controller can
accept a frame and acknowledge its checksum without rotating the shaft. Confirm the
mode before concluding that a ``no motion'' result is a wiring or CAN problem.
\end{caution}

\section{The MKS CAN protocol}

\subsection{Frame structure and checksum}

Every frame is \texttt{[command byte][data bytes][checksum byte]}, where the checksum
is an 8-bit sum that \textbf{includes the CAN ID}:

\begin{lstlisting}
def checksum(motor_id, data_bytes):
    return (motor_id + sum(data_bytes)) & 0xFF
\end{lstlisting}

\begin{center}
\begin{tabular}{llll}
\toprule
\textbf{Frame, no checksum} & \textbf{CAN ID} & \textbf{Arithmetic} & \textbf{Checksum} \\ \midrule
\texttt{31}    & \texttt{0x04} & \texttt{0x31+0x04}       & \texttt{35} \\
\texttt{31}    & \texttt{0x06} & \texttt{0x31+0x06}       & \texttt{37} \\
\texttt{3A}    & \texttt{0x05} & \texttt{0x3A+0x05}       & \texttt{3F} \\
\texttt{F3 01} & \texttt{0x05} & \texttt{0xF3+0x01+0x05}  & \texttt{F9} \\
\texttt{F7}    & \texttt{0x06} & \texttt{0xF7+0x06}       & \texttt{FD} \\
\bottomrule
\end{tabular}
\end{center}

The checksum differs per joint, so a frame copied from another joint's documentation
is silently rejected. \texttt{BASIC\_CAN\_COMMANDS} has a ready-made table for every
ID on this arm.

\begin{caution}{Never hand-invent a checksum, and never reuse another joint's frame}
Omitted or incorrect checksum bytes were the direct cause of several ``the motor
ignored my command'' failures in this project. Compute it from the formula every
time, including for the emergency stop.
\end{caution}

\subsection{Command reference}

\begin{longtable}{lll}
\toprule
\textbf{Name} & \textbf{Byte} & \textbf{Status} \\ \midrule
\endhead
\texttt{QUERY\_MOTOR}       & \texttt{0xF1} & Liveness check. Answers \texttt{F1 01}. \\
\texttt{READ\_ENCODER}      & \texttt{0x31} & \textbf{Verified. Use this one.} \\
\texttt{READ\_ENCODER} carry form & \texttt{0x30} & Superseded. See Section~\ref{sec:corrections}. \\
\texttt{READ\_VELOCITY}     & \texttt{0x32} & Responds; zero at rest. Layout unconfirmed. \\
\texttt{READ\_ERROR}        & \texttt{0x39} & Signed 32-bit deficit. Grows when stalled. \\
\texttt{READ\_ENABLE\_STATE} & \texttt{0x3A} & \texttt{01} enabled, \texttt{00} disabled. Verified. \\
\texttt{READ\_SHAFT\_PROTECTION} & \texttt{0x3E} & \texttt{00} clear. Non-zero latches the driver. \\
\texttt{READ\_IO}             & \texttt{0x34} & Input and output pin states. Verified. Section~\ref{sec:endstops}. \\
\texttt{SET\_WORKING\_MODE} & \texttt{0x82} & See the panel caution. \\
\texttt{SET\_HOME\_PARAMS}   & \texttt{0x90} & Trigger level, direction, speed, \texttt{EndLimit}. \\
\texttt{GO\_HOME}            & \texttt{0x91} & Runs onto the limit. \textbf{Unlocks the shaft} if \texttt{EndLimit=1}. \\
\texttt{LIMIT\_PORT\_REMAP}  & \texttt{0x9E} & \textbf{Required for 42D endstops here.} No read-back. \\
\texttt{ENABLE\_SHAFT\_PROTECTION} & \texttt{0x88} & \textbf{Avoid.} False positives, and it latches. \\
\texttt{ENABLE\_MOTOR}      & \texttt{0xF3} & \texttt{01} enable, \texttt{00} disable. Verified. \\
\texttt{RELATIVE\_POSITION} & \texttt{0xF4} & Unsafe packing. See the motion section. \\
\texttt{ABSOLUTE\_POSITION} & \texttt{0xF5} & \textbf{Unsafe.} See the motion section. \\
\texttt{SPEED\_CONTROL}     & \texttt{0xF6} & Not exercised. \\
\texttt{EMERGENCY\_STOP}    & \texttt{0xF7} & Verified, with checksum. \\
\texttt{POSITION\_CONTROL}  & \texttt{0xFD} & \textbf{Verified. Preferred for all motion.} \\
\bottomrule
\end{longtable}

\subsection{Reading the encoder}
\label{sec:encoder}

Use \texttt{0x31}. The response is \texttt{[0x31][6 bytes signed position][checksum]},
where the six bytes are a big-endian signed 48-bit multi-turn count on the
16384-counts-per-revolution scale:

\begin{lstlisting}
raw48 = int.from_bytes(data[1:7], byteorder="big", signed=True)
motor_degrees = raw48 * 360.0 / 16384
joint_degrees = motor_degrees / gear_ratio
\end{lstlisting}

This is multi-turn and does not wrap at one revolution. It resets to approximately
zero when the controller is power-cycled, so it is a relative reference from
power-on, not an absolute machine coordinate.

\begin{caution}{The encoder measures the motor shaft, not the joint output}
Gearbox backlash and any lost motion downstream of the motor are invisible to this
encoder. It cannot measure output-side backlash; that needs a dial indicator on the
output.
\end{caution}

\subsection{Encoder non-linearity}

The magnetic encoder carries a once-per-revolution sinusoidal error caused by the
sensing magnet sitting slightly off-centre on the shaft. Measured on Joint A over
6725 samples: \textbf{1.465\textdegree{} peak to peak at the motor shaft}, about
0.407\% of a revolution.

Two practical consequences:

\begin{itemize}[leftmargin=*]
    \item It cancels over any whole number of motor revolutions. Choosing a step size
          that is an integer number of motor turns makes it vanish --- Joint A steps
          of exactly 6 motor revolutions landed at 100.0\% every time, where
          90\textdegree{} motor-shaft steps scattered between 98.7 and 101.4\%.
    \item Through a reduction it becomes negligible at the joint: 1.465\textdegree{}
          through 48:1 is 0.031\textdegree{}.
\end{itemize}

Repeatability is unaffected. Returning to the same shaft angle repeats to better
than one encoder count. Absolute readings carry the cyclic error; relative and
repeat positioning do not.

\section{Motion commands}
\label{sec:motion}

\textbf{Use \texttt{POSITION\_CONTROL}, \texttt{0xFD}.} It is what the project's own
working jog code uses and the only motion command verified safe here.

\begin{lstlisting}
[0xFD]
[direction | ((speed >> 8) & 0x0F)]   direction 0x00 fwd, 0x80 reverse
[speed & 0xFF]                        speed is approximately motor rpm
[accel]
[(pulses >> 16) & 0xFF]               ALWAYS-POSITIVE magnitude, on the
[(pulses >>  8) & 0xFF]               3200-pulses-per-motor-rev scale
[pulses & 0xFF]
[checksum]
\end{lstlisting}

\[
\text{pulses} = |\text{joint degrees}| \times \frac{3200 \times \text{gear ratio}}{360}
\]

The driver replies with an unsolicited two-frame status sequence: \texttt{FD 01} for
``started'', then \texttt{FD 02} for ``complete'' if the commanded distance was
actually achieved. \textbf{The absence of \texttt{FD 02} is the reliable signal that
a command did not finish.}

\begin{caution}{Do not use \texttt{0xF5}, and be careful with \texttt{0xF4}}
\texttt{0xF5} packs an absolute target through a plain \texttt{\& 0xFFFFFF}. On a
joint whose raw counter had drifted to roughly $-274000$ counts, a small intended
nudge became an enormous, non-settling, uninterruptible move that required cutting
motor power. \texttt{0xF4} packs a possibly-negative delta the same way. The
\texttt{0xFD} pattern above, a separate direction byte plus an always-positive
magnitude, does not have this failure mode.
\end{caution}

\begin{caution}{Never retry a stalled command into an obstruction}
Re-sending the same command to a joint that has not moved changes nothing except the
motor's temperature. Nine consecutive full-torque retries into a hard stop drove one
joint's internal position error to $-124374$ counts and cost a cooldown to recover
from. If a step under-delivers, stop, or back off and re-attempt once.
\end{caution}

\begin{caution}{A repeatable stop angle is mechanical, not roughness}
Rough gear mesh varies. If a joint stops at the \emph{same} angle to the encoder
count on attempt after attempt, and the back-off also delivers nothing, that is a
hard mechanical limit. On Joint B this pattern turned out to be a gear seated in the
wrong position; once corrected the joint swept the same range with zero unstick
cycles. Open the gearbox rather than working at it.
\end{caution}

\section{Speed and its limits}

The \texttt{speed} field is very close to motor rpm. Measured on a bare NEMA 17 with
no load:

\begin{center}
\begin{tabular}{lll}
\toprule
\textbf{Commanded} & \textbf{Actual} & \textbf{Behaviour} \\ \midrule
10--225  & tracks to within 0.5\%      & linear \\
250      & 98.7\%                      & edge of linear \\
275--325 & 96.5\% down to 90\%         & rolling off \\
350 plus & saturates near 310 rpm      & losing steps \\
\bottomrule
\end{tabular}
\end{center}

\subsection{Measured ceilings, per joint}

That table is a bare motor on a bench. A joint drags its gearbox with it and the
ceiling falls accordingly, so treat the no-load numbers as an upper bound and
measure each joint.

\begin{center}
\begin{tabular}{lllll}
\toprule
\textbf{Joint} & \textbf{Driver} & \textbf{Gear} & \textbf{Linear to} & \textbf{Saturates} \\ \midrule
bench, no load & SERVO42D & none    & 225 rpm & about 310 rpm \\
B              & SERVO42D & 67.82:1 & 150 rpm & about 165 rpm \\
\bottomrule
\end{tabular}
\end{center}

Joint B was measured on 2026-09-02 over a 48\textdegree{} arc. Commanded 150 tracked
at 99.7\%; commanded 165, 175, 200, 225 and 300 all produced 159 to 165 rpm. Past
saturation, extra commanded speed buys nothing but current and heat. The measurement
was repeated at eight times the acceleration to confirm the ramp was not the limiter,
which moved the result by only about 4\%.

\begin{caution}{Over-commanding speed did not lose position on a geared joint}
\texttt{0xFD} commands a pulse \emph{count}, not a duration. On Joint B every trial
from 150 through 300 commanded delivered its 48\textdegree{} arc to $+0.01\%$: past
saturation the move simply takes longer, it does not drop pulses. Silent position
loss requires actual slip. Do not carry the no-load ``loses steps above 250'' result
across to a geared joint automatically --- verify which behaviour the joint in front
of you shows.
\end{caution}

\begin{caution}{Heat is the real limit on any speed or reliability testing}
A ladder of trials run back to back overheats the motor even when no single trial is
aggressive. Eleven trials with one-second gaps was enough to have to stop and let a
joint cool; fifteen seconds between trials was not. There is no temperature-read
opcode on these drivers, so software cannot detect this. Budget total energized time,
not just peak current, and stop as soon as you have the answer.
\end{caution}

\section{Converting encoder counts to joint degrees}

\[
\text{counts per joint degree} = \frac{\text{gear ratio} \times 16384}{360}
\]

\begin{center}
\begin{tabular}{lll}
\toprule
\textbf{Gear ratio} & \textbf{Counts per joint degree} & \textbf{Joints} \\ \midrule
13.5:1  & 614.4   & X \\
48.0:1  & 2184.5  & A \\
67.82:1 & 3086.56 & B, C \\
150.0:1 & 6826.7  & Y, Z \\
\bottomrule
\end{tabular}
\end{center}

\begin{caution}{Three wrong conversion constants exist in this project's history}
Do not use any of them.
\begin{enumerate}[leftmargin=*]
    \item ``About 3300 counts per 90\textdegree'' describes the \emph{motor shaft},
          not the joint. Through 67.82:1 that same motion is only about
          1.07\textdegree{} of joint travel.
    \item \texttt{COUNTS\_PER\_REV = 16050}. The encoder is 14-bit, so 16384 is
          correct and source-verified.
    \item \texttt{pulses\_per\_degree = -813.0} in \texttt{mks\_driver.py}, off by
          roughly 3.8 times, with an accompanying ``safety cage'' that estimated
          limits using a different multiplier than the one actually transmitted, so
          it could not reliably block an unsafe command.
\end{enumerate}
Use only the formula above.
\end{caution}

\section{Endstops and homing}
\label{sec:endstops}

\subsection{Where the sensor actually connects}
There are two wiring schemes and they are mutually exclusive. The vendor default
puts a limit switch on \texttt{EVCC}, \texttt{EGND} and \texttt{IN\_1}, all on the
same green terminal block as the CAN pair --- top to bottom, that block is
\texttt{5V}, \texttt{GND}, \texttt{IN}, \texttt{CAN\_H}, \texttt{CAN\_L}. The
silkscreen is split, \texttt{5V/GND/IN} printed above the block and \texttt{CAN H L}
below it, which is easy to misread if you only look at one end.

\textbf{This build does not use that scheme.} It follows the Arctos documentation,
which for the 42D axes enables \emph{limit port remap} and moves the limits to the
opposite connector: left limit to \texttt{En}, right limit to \texttt{Dir}, with
\texttt{Com} tied to the matching high level. \texttt{En} and \texttt{Dir} are
opto-isolated, so with \texttt{Com} floating they register nothing at all.

\begin{caution}{The 42D has no \texttt{IN\_2}, which is why the remap exists}
Per the vendor manual's IO table, \texttt{IN\_1} (home / left limit) is present on
both 57D and 42D, but \texttt{IN\_2} (right limit) is 57D only. Remapping is how a
42D gets a second endstop. On a 57D the remap is optional, since that board has a
real \texttt{IN\_2}.
\end{caution}

\begin{caution}{No DIP switch is involved on the 42D}
The \texttt{SW1} pin-2 instruction that circulates for these boards belongs to the
57D. On the 42D, \texttt{EVCC}/\texttt{EGND} is the board's own 5\,V, rated 20\,mA
--- the whole budget for the sensor and its indicator LEDs.
\end{caution}

\subsection{Enabling the remap}
\begin{lstlisting}
# LIMIT_PORT_REMAP (0x9E), enable = 01, CAN ID 0x06
# checksum 0x06 + 0x9E + 0x01 = 0xA5
cansend can0 006#9E01A5
# reply 9E 01 A5 -> status 1, set success
\end{lstlisting}

The manual documents \texttt{0x9E} as a write only. \textbf{Nothing reports whether
remap is currently on}, so it cannot be queried --- only set. Two further conditions
must hold or it does nothing useful: the driver must be in a serial
(\texttt{SR\_}) work mode, because in any \texttt{CR\_} mode \texttt{En} and
\texttt{Dir} are the pulse interface's own enable and direction inputs; and
\texttt{Com} must be wired.

\subsection{Reading the pins: \texttt{0x34}}
The reply is \texttt{[0x34][status][checksum]}. The status byte:

\begin{center}
\begin{tabular}{lll}
\toprule
\textbf{Bit} & \textbf{Pin} & \textbf{After remap, reports} \\ \midrule
bit 0 & \texttt{IN\_1} & the \texttt{En} pin \\
bit 1 & \texttt{IN\_2} & the \texttt{Dir} pin \\
bit 2 & \texttt{OUT\_1} & stall indication, 57D only \\
bit 3 & \texttt{OUT\_2} & reserved, 57D only \\
bits 4--7 & \multicolumn{2}{l}{reserved} \\
\bottomrule
\end{tabular}
\end{center}

Reading \texttt{0x34} does \emph{not} require \texttt{EndLimit}; that setting only
governs whether the driver acts on a limit. Measured on Joint C with a dual-output
hall sensor, one output per magnet pole:

\begin{center}
\begin{tabular}{lll}
\toprule
\textbf{Byte} & \textbf{Meaning} & \\ \midrule
\texttt{0x03} & both idle & resting state \\
\texttt{0x02} & bit 0 low & north pole on \texttt{En} \\
\texttt{0x01} & bit 1 low & south pole on \texttt{Dir} \\
\bottomrule
\end{tabular}
\end{center}

Both are \textbf{active low}, so the default \texttt{Hm\_Trig = Low} is correct and
needs no change.

\begin{caution}{A resting \texttt{0x01} before remap looks identical to a dead sensor}
Without the remap, bit 0 reports the physical \texttt{IN\_1} terminal. On this build
that terminal is deliberately empty, so it idles high forever regardless of the
magnet. Two sessions and two different sensors were spent on this before the cause
was found. If an endstop reads nothing, re-send \texttt{9E 01} before suspecting the
wiring.
\end{caution}

\subsection{Homing}
\texttt{SET\_HOME\_PARAMS (0x90)} carries \texttt{homeTrig} (0 low, 1 high),
\texttt{homeDir} (0 CW, 1 CCW), a 16-bit \texttt{homeSpeed} in rpm, and
\texttt{EndLimit} (0 disable, 1 enable). \texttt{GO\_HOME (0x91)} answers
\texttt{0} fail, \texttt{1} started, \texttt{2} complete. The manual requires a
\texttt{GoHome} after enabling \texttt{EndLimit} or changing any limit parameter.

\begin{caution}{\texttt{EndLimit=1} unlocks the shaft at the left limit}
The manual is explicit: when \texttt{EndLimit} is enabled and a \texttt{GoHome}
reaches the left limit switch, the driver \textbf{unlocks the shaft}. On a belted,
gravity-loaded joint that is a drop. Test with the belt off. Set
\texttt{EndLimit=0} if the axis must stay locked after homing.
\end{caution}

\section{Safety}
\label{sec:safety}

\begin{center}
\begin{tabular}{lll}
\toprule
\textbf{Action} & \textbf{Frame} & \textbf{Example, Joint A at \texttt{0x04}} \\ \midrule
Enable coils   & \texttt{F3 01 crc} & \texttt{cansend can0 004\#F301F8} \\
Disable coils  & \texttt{F3 00 crc} & \texttt{cansend can0 004\#F300F7} \\
Emergency stop & \texttt{F7 crc}    & \texttt{cansend can0 004\#F7FB} \\
\bottomrule
\end{tabular}
\end{center}

Recompute for the joint you are actually driving, or take the row from
\texttt{BASIC\_CAN\_COMMANDS}.

\begin{caution}{Compute your joint's e-stop frame before you need it}
A previous version of this documentation gave a joint's emergency stop as
\texttt{cansend can0 001\#F7F8}, addressed to CAN ID \texttt{0x01}, when the motor
was not on that ID. It would have done nothing. Derive the frame for the joint in
front of you, and keep it pasted in a second terminal before commanding motion.
\end{caution}

\begin{caution}{Dropping coils on a gravity-loaded joint lets it fall}
Several scripts in this package disable the motor in a \texttt{finally} block. On an
unloaded bench motor that is harmless; on an assembled arm it removes holding
torque. Know which you have before running one.
\end{caution}

\section{Troubleshooting}
\label{sec:troubleshooting}

\begin{longtable}{p{0.34\textwidth}p{0.58\textwidth}}
\toprule
\textbf{Symptom} & \textbf{Likely cause and fix} \\ \midrule
\endhead
\texttt{can0} does not exist &
  \texttt{slcand} died, usually because the adapter re-enumerated. Re-run the
  bring-up. Check \texttt{/dev/serial/by-id/} for the current device name. \\
Adapter present, \texttt{can0} missing after a replug &
  Expected. \texttt{slcand} does not restart itself. A udev rule keyed on the
  adapter's serial can automate this. \\
\texttt{ip} shows \texttt{bitrate 0} &
  Normal on slcan. Not a fault. \\
No response from any ID &
  Check panel CAN ID, CAN rate and CANRSP; check CAN\_H, CAN\_L and termination;
  confirm the driver is powered. If still silent see Section~\ref{sec:rs485}. \\
Command acknowledged, motor does not move &
  Distinguish ``not being driven'' from ``driven but cannot move the load'' --- the
  bus cannot tell these apart. Disable coils and turn the joint by hand. \\
Motion for a few degrees, then nothing in either direction &
  If it cannot even back off, suspect a bind stiff enough to absorb full torque, a
  marginal motor phase connection, or a misassembled gear. Reseating a connector
  once took a joint's travel budget from 1.05\textdegree{} to 3.85\textdegree{}. \\
Stops at the same angle every attempt &
  Mechanical, not roughness. Open the gearbox. A gear in the wrong position produced
  exactly this on Joint B. \\
Gritty resistance by hand, stalls under power &
  Gear mesh needs lubrication. Use a plastic-safe PTFE or silicone grease.
  Petroleum jelly works as a diagnostic but migrates once warm. \\
Motor warm but not hot &
  Normal. Thermal shutdown makes a motor too hot to touch; warm is a stepper doing
  its job. \\
Driver stops responding mid-test &
  Check power and connectors first. This has been an accidental unplug more than
  once. \\
\textbf{Encoder tracks perfectly but the arm does not move} &
  \textbf{Suspect the drivetrain, not the driver.} The encoder is on the motor
  shaft, upstream of the reduction, so slipping cycloidal gears produce a flawless
  trace --- delivered counts matching commanded --- for an output going nowhere.
  Slip also \emph{unloads} the motor, so no stall, torque or position-error
  heuristic will fire. Confirm motion at the output before trusting any encoder
  result. \\
Endstop never changes \texttt{0x34}, sensor LED does light &
  On a 42D, limit port remap is probably not enabled, so bit 0 is reporting the
  empty \texttt{IN\_1} terminal rather than \texttt{En}. Send \texttt{9E 01}.
  There is no read-back, so it cannot be confirmed by query.
  Section~\ref{sec:endstops}. \\
Both endstop inputs read low at rest &
  Not ``nothing connected'' --- an unwired input reads \emph{high}. Most likely a
  reversed 3-pin sensor plug: \texttt{+V} sits in the middle, so reversing it
  swaps \texttt{Signal} and \texttt{GND} and straps the driver input to the
  module's ground while leaving the sensor powered and its LED lit. Turn the plug
  around. \\
Sensor LED lights but the \texttt{0x34} bit does not move &
  The sensor is detecting; the signal is not reaching the input pin. The fault is
  in the wire between them. Watching LED and bit together is what localises this. \\
Endstop bit flickers rapidly while parked &
  A magnet sitting at the edge of the sensor's range, crossing its threshold
  repeatedly. Release edges are narrow --- 0.062\textdegree{} measured on Joint X.
  Move the sensor or magnet; do not fix it in software. \\
A move continues after the script exits &
  Expected. \texttt{ENABLE\_MOTOR 00} does not cancel an in-flight \texttt{0xFD};
  only \texttt{EMERGENCY\_STOP (0xF7)} does. Send \texttt{0xF7} before the coil
  disable. \\
Motors hold torque the moment they are powered &
  The \texttt{En} option is set to \texttt{Hold} (always enabled). Set it back to
  active-low with \texttt{SET\_EN\_ACTIVE (0x85)} data \texttt{00}. Do not apply
  this to a joint using limit port remap, where \texttt{En} is a limit input. \\
\bottomrule
\end{longtable}

\subsection{Total silence with a powered driver: the RS485 variant}
\label{sec:rs485}

The MKS SERVO42D and SERVO57D ship in both CAN and RS485 variants. The RS485 boards
have \textbf{no CAN transceiver at all} and will never appear on the bus at any
bitrate or address. Joint X's driver turned out to be the RS485 variant, which is why
it has never communicated. If a joint is silent across a full ID scan and a bitrate
sweep, with power and wiring confirmed, \textbf{check the board's part number before
spending any more time in software}.

\section{The scripts in this package}

All of these live in \texttt{\textasciitilde/arctos\_ws/src/arctos\_can\_control/arctos\_can\_control}.
Each takes \texttt{--can-id} and \texttt{--gear-ratio}; the motion ones also take an
explicit safety band.

\begin{longtable}{p{0.33\textwidth}p{0.59\textwidth}}
\toprule
\textbf{Script} & \textbf{What it does} \\ \midrule
\endhead
\texttt{can\_id\_scan.py} &
  Read-only sweep of CAN IDs using \texttt{0x31}. Never enables coils. Run this
  before trusting any ID. \\
\texttt{joint\_reliability\_test.py} &
  Statistical encoder sampling at rest. Response rate, checksum failures, jitter,
  drift. \\
\texttt{joint\_unstick\_move\_test.py} &
  Stepped \texttt{0xFD} travel in one direction with a back-off-and-retry cycle for
  a gritty mesh. \\
\texttt{joint\_sweep\_unstick\_test.py} &
  Sweeps back and forth between two bounds for a fixed wall-clock duration, with the
  same unstick cycle and a total stall budget as a heat guard. Use when a rough spot
  needs crossing repeatedly. \\
\texttt{joint\_speed\_limit\_test.py} &
  Ladder of commanded speeds across a fixed arc. Reports steady rpm, a whole-move
  average as a cross-check, and delivered position against commanded. \\
\texttt{joint\_position\_control\_test.py} &
  Stepped position moves with retry. Superseded by the unstick variants for a
  suspect mesh. \\
\texttt{joint\_endstop\_test.py} &
  Read-only live monitor of the \texttt{0x34} byte. Sends nothing else, so it is safe
  on an assembled arm. Prints only when the byte changes. \\
\texttt{joint\_endstop\_dwell.py} &
  Timestamps every endstop edge and reports how long each trigger held, with a
  solid/chatter/marginal verdict. Use it when the monitor above shows a trigger but
  you need to know whether it is stable; that script has no timestamps, so a 50\,ms
  tap and a two-second hold look identical there. \\
\texttt{joint\_io\_diff.py} &
  Sweeps every documented read-only register with the magnet away and then held on,
  reporting any byte that differs. For when an endstop does not appear on
  \texttt{0x34} at all. \\
\texttt{joint\_magnet\_probe.py} &
  Checks whether an external magnet disturbs the encoder reading. \\
\bottomrule
\end{longtable}

\section{ROS 2 package layout}

The workspace is a standard ROS 2 layout. The CAN work lives in
\texttt{src/arctos\_can\_control}; \texttt{src/ros2\_arctos} carries the description,
MoveIt configuration and hardware interface.

\begin{lstlisting}
cd ~/arctos_ws
colcon build --packages-select arctos_can_control
source install/setup.bash
\end{lstlisting}

The Python test scripts above do not need the workspace built or sourced. They talk
to SocketCAN directly and can be run from the source directory, which is why every
procedure in the SOPs uses them rather than ROS 2 nodes.

\section{Corrections, and what this set replaces}
\label{sec:corrections}

These four documents replace five older ones. The following items were
\textbf{wrong} in those files and were corrected here, so do not re-derive them from
any surviving copy.

\begin{itemize}[leftmargin=*]
    \item \textbf{Joint B and C CAN IDs.} Joint B is \texttt{0x05} and Joint C is
          \texttt{0x06}, as the older files and the controller config state. An
          interim revision of this document asserted the reverse after a scan found
          Joint B answering on \texttt{0x06}; that driver's panel ID was wrong, and
          the map was right. Scan before driving, and correct the panel.
    \item \textbf{Encoder opcode.} An old guide standardized on \texttt{0x30} and
          described \texttt{0x31} as unverified. That is backwards. \texttt{0x31} is
          what the real C++ driver uses and what every script here uses.
    \item \textbf{Bitrate verification.} An old guide instructed the reader to run
          \texttt{ip link set can0 up type can bitrate 500000} and then verify by
          looking for \texttt{bitrate 500000}. Neither works on slcan.
    \item \textbf{Bus error counters.} Old troubleshooting advice leaned on bus-error
          counts that the slcan driver never populates.
    \item \textbf{Motion command.} An old safe-first-motion procedure was built on
          \texttt{0xF5}, whose absolute-versus-relative behaviour it correctly
          flagged as unresolved. That question is moot: use \texttt{0xFD}.
    \item \textbf{Missing checksums.} The old cheat sheet listed
          \texttt{cansend can0 005\#F300} and \texttt{cansend can0 005\#F7} with no
          checksum byte at all. Both would be rejected. The correct frames are
          \texttt{005\#F300F8} and \texttt{005\#F7FC}.
    \item \textbf{Speed ceiling.} The no-load figure of 225 rpm was presented as a
          general operating limit. It is a bare-motor number; Joint B saturates at
          about 165.
\end{itemize}

\end{document}
